\subsection{The hyperelliptic refinement of the scalar obstruction}

\begin{claim}{AC.41}{Proved; an all-extension obstruction using the pointed involution}
Let $C/k$ be a genus-two curve over a finite field of
characteristic two, and let $P\in C(k)$ be a Weierstrass
point. Fix a nontrivial geometric connected cyclic etale
cover of odd degree $n$, with a chosen point above $P$.
Let $k_d/k$ be the first constant extension over which
its geometric isomorphism class is stable, and use its
unique pointed descent to $k_d$. Suppose Frobenius on
its geometric cyclic deck group is multiplication by
$a\in(\mathbf Z/n)^\times$, of order $m$. Put
\[
 U=(\mathbf Z/n)^\times,\qquad
 T=\langle2\rangle\subset U,\qquad
 H=\langle a,-1\rangle\subset U.
\]
If $T\subset H$, every geometric two-torsion point of
the primitive factor is rational over the first
constant-deck field $k_{dm}$. Consequently, for this
pointed cover, the primitive class number is odd over
some constant-deck field if and only if its geometric
primitive two-rank is zero. When $n=\ell$ is prime,
this is the whole relative criterion \textup{(R)}.

More generally, the possible primitive Frobenius
multipliers over $k_{dm}$ belong to the units of a
direct factor of
\[
 \left(\mathbf F_{2^{\delta}}\right)^c,
 \qquad
 \delta=\frac{|T|}{|T\cap H|},\qquad
 c=\frac{|U|}{|TH|}.                              \tag{AC.41.1}
\]
Thus the preceding sufficient condition is exactly
$\delta=1$ for this fixed-ring argument.

In particular, suppose the Jacobian Frobenius of $C/k$
is the scalar $[\varepsilon2^r]$, with
$\varepsilon\in\{1,-1\}$ and $\#k=2^{2r}$.
For an odd prime $\ell$, put $f=\operatorname{ord}_\ell(2)$.
The sufficient condition becomes, independently of the
sign $\varepsilon$,
\[
 \begin{cases}
  \gcd(r,f)=1,&f\text{ odd},\\
  \gcd(r,f/2)=1,&f\text{ even}.
 \end{cases}                                      \tag{AC.41.2}
\]
This sharpens AC.38, and also applies to nonscalar
Jacobian Frobenius whenever the displayed deck-group
condition holds over the first field of definition of
the chosen cyclic cover.
\end{claim}
\begin{proof}
The hyperelliptic involution $\iota$ fixes $P$ and acts
as minus one on $J_C$. It therefore preserves the
isomorphism class of every cyclic cover. It has a
unique lift $\tau$ fixing the chosen point above $P$:
two such lifts differ by a deck transformation fixing
that point, hence are equal. The lift inverts the
cyclic deck group, and its square is the identity.
Uniqueness also shows that the lift is defined over
$k_d$ and commutes with its arithmetic Frobenius $F$.
This uses the pointed descent; an arbitrary unpointed
twist need not preserve this commutation relation.

Let
\[
 R=\mathbf F_2[X]/\Phi_n(X),\qquad
 W=H^1_{\mathrm{et}}(D_{\overline k},\mathbf F_2)_{\mathrm{prim}}.
\]
Since $n$ is odd, $R$ is finite etale. Every nontrivial
coherent character space of a genus-two cyclic cover
has dimension one. The geometric etale part therefore
has rank zero or one over every field factor of $R$;
equivalently, $W$ is free of rank one over a direct
factor $R'$ of $R$. This is the same multiplicity
argument used in AC.38. The support $R'$ is stable
under $X\mapsto X^a$ and $X\mapsto X^{-1}$, because
$F$ and $\tau$ act invertibly on $W$.

The action of $F$ is semilinear for the deck
automorphism $X\mapsto X^a$ (or its inverse under the
dual cohomology convention). Since $a^m=1$, $F^m$ is
$R'$-linear, and hence is multiplication by a unit
$b\in R'$. Commuting $F^m$ with $F$ shows that
$b$ is fixed by $a$. Commuting $F^m$ with $\tau$
shows that $b$ is fixed by inversion. Thus
\[
                         b\in(R'^H)^\times.           \tag{AC.41.3}
\]
This argument does not assume that the primitive
factor is ordinary; rank zero factors have simply
been omitted from $R'$.

Over $\overline{\mathbf F}_2$, the primitive character
set is $U$. The field factors of $R$ are the orbits
of multiplication by two, that is, the cosets of $T$.
Taking $H$-invariants groups these into cosets of $TH$.
The stabilizer $T\cap H$ acts on each corresponding
field of degree $|T|$, leaving a field of degree
$|T|/|T\cap H|$. This proves (AC.41.1). A stable
direct factor $R'$ merely removes some of these
field factors.

If $T\subset H$, the fixed ring is a product of copies
of $\mathbf F_2$, whose unit group is trivial. Hence
$b=1$ and $F^m$ is the identity on $W$. The same holds
on the dual geometric two-torsion points of the
primitive abelian factor. A nonzero such point then
remains rational over every further constant-deck
extension. Conversely, geometric two-rank zero
excludes rational two-torsion over every field, so
its class number is always odd. Any field over which
this pointed geometric cover has constant decks
contains $k_{dm}$: first the isomorphism class must
be stable, requiring a multiple of $d$, and then
the deck scalar must become the identity, requiring
a further multiple of $m$.

For scalar Jacobian Frobenius $[s]$ with
$s=\varepsilon2^r$ and $s^2=\#k$, every cyclic
isomorphism class is stable over $k$ and the pointed
deck scalar is $s$. Equivalently, on its Kummer line
the scalar is $s$, and Cartier duality gives the
deck scalar $(\#k)/s=s$. Therefore
$H=\langle2^r,-1\rangle$. For prime $\ell$, if
$f$ is odd then $-1\notin T$, and
$H\cap T=\langle2^r\rangle$; the fixed-field degree
is $\gcd(r,f)$. If $f$ is even, $-1=2^{f/2}$ in
$T$, so $H=\langle2^{\gcd(r,f/2)}\rangle$ and
the fixed-field degree is $\gcd(r,f/2)$. This proves
(AC.41.2).

For composite $n$ the same inclusion and fixed-ring
formulas remain valid. The simple prime formula
must then be replaced by the actual subgroup
calculation in $(\mathbf Z/n)^\times$; in particular,
an even order of two does not by itself imply
$-1\in\langle2\rangle$.
\end{proof}

\paragraph{Scope for the counterexample pencil.}
For a member of the supersingular pencil over
$\mathbf F_{2^{38}}$, the scalar is $-2^{19}$.
If it has no nontrivial cyclic primitive-zero cover
as in AG.33--AG.34, (AC.41.2) excludes every rational
prime-cyclic \textup{(R)} witness whose prime $\ell$
satisfies $19\nmid\operatorname{ord}_\ell(2)$,
over every finite extension. In particular, every
such witness has $\ell\ge191$: for $\ell<191$,
divisibility by19 would force $\ell=19j+1$ with
$j\in\{2,4,6,8\}$, and the four values
$39,77,115,153$ are composite. This is a necessary
condition on witnesses, not a claim that degree191
succeeds or that the arithmetic criterion A3 fails.
